\ParallelTitle{放慢脚步，重看熟悉的街道}{いつもの通りを、ゆっくり歩く}

\ParallelSection{notice}{留一点时间观察}{気づくための時間をつくる}

\ParallelText{opening}{一条熟悉的街道，有时只是两项日程之间的一段路。不妨试着走一次，不急着完成任何差事。目的不是寻找引人注目的地标，也不是拍出完美的照片，而是留意那些在匆忙赶路时被忽略的细节。}{いつもの通りが、予定と予定をつなぐだけの道になっていることがあります。一度、用事を済ませるためではなく歩いてみましょう。立派な名所を探したり、完璧な写真を撮ったりするのではなく、急いでいると見落とす細部に目を向けるのが目的です。}

\ParallelText{change}{一扇刷过漆的门、一片树荫，或敞开的窗户里传来的碗碟声，都可以成为观察的起点。无须立刻解释每一个细节。有价值的观察可能先是一个问题；散步时可以暂时保留它，等有空再进一步了解。}{塗り直された扉、木陰、開いた窓から聞こえる食器の音。それだけでも観察のきっかけになります。すべてをその場で説明する必要はありません。問いから始まる気づきもあります。調べる時間ができるまで、その問いを残しておいてもよいのです。}

\ParallelSection{route}{选一条简单的路线}{無理のない道順を選ぶ}

\ParallelText{route-plan}{起点最好方便到达，终点也不应让你必须匆忙折返。比起排得满满的行程，一小段环线往往更合适。出行时仍需查看当地的实际情况；下图是虚构的示意图，并不是某个真实地点的路线指引。}{行きやすい出発点と、帰りを急がずに済む終点を選びます。欲張った行程より、短い周回路のほうが役立つこともあります。歩く場所の状況は実際に確かめてください。下の図は架空の例で、実在する場所への道案内ではありません。}

\ParallelFigure{route-image}{images/figure-route-image.png}{虚构路线上的四个停留点。数字只标识停留位置，不代表距离、时间或无障碍条件。}{四つの立ち止まる場所を示した架空の道順です。数字は場所の目印であり、距離、所要時間、バリアフリー状況を示すものではありません。}

\ParallelText{steps.item-1}{\begin{itemize}\item \phantomsection\label{steps}\hypertarget{steps}{}每次停留只选一个细节，不必试图描述周围的一切。\end{itemize}}{\begin{itemize}\item \ifdefstring{\ParallelMode}{right}{\phantomsection\label{steps}\hypertarget{steps}{}}{}立ち止まるたびに一つの細部を選び、周囲のすべてを説明しようとしないこと。\end{itemize}}

\ParallelText{steps.item-2}{\begin{itemize}\item 留一点没有安排的时间。遇到值得再看一眼的事物，可以多停留一会儿。\end{itemize}}{\begin{itemize}\item 予定を入れない時間も残しましょう。もう一度見たいものがあれば、その場にいられます。\end{itemize}}

\ParallelText{steps.item-3}{\begin{itemize}\item 留意他人的隐私。写下一项观察，并不需要给陌生人拍照。\end{itemize}}{\begin{itemize}\item 周囲の人のプライバシーに配慮しましょう。観察を書き留めるために、見知らぬ人を撮影する必要はありません。\end{itemize}}

\clearpage

\ParallelSection{notebook}{带一本小笔记本}{小さなノートに書き留める}

\ParallelText{notebook-prose}{笔记本里可以写片段，不必句句都完整漂亮。先记录实际看到的内容，再补充解释。“长椅是空的”是一项观察；“没有人来这个广场”则是范围大得多的判断。把这两类表述分开，之后写出来的文字会更诚实，也更有意思。}{ノートには、整った文章ではなく断片を書いてもかまいません。解釈を加える前に、実際に気づいたことを記録します。「ベンチは空いていた」は観察ですが、「この広場は誰も使わない」はずっと広い主張です。二つを区別すれば、後で書く文章がより誠実で興味深いものになります。}

\ParallelText{prompt}{\begin{quote}一种声音、一种颜色、一个问题：给自己的笔记一个简单的提示。\end{quote}}{\begin{quote}音を一つ、色を一つ、問いを一つ。自分のメモのための簡単な手がかりです。\end{quote}}

\begin{ParallelKeep}

\ParallelText{notes-table.row-0}{\phantomsection\label{notes-table}\hypertarget{notes-table}{}\begin{tabularx}{\linewidth}{@{}>{\RaggedRight\arraybackslash}X>{\RaggedRight\arraybackslash}X@{}}\toprule \textbf{观察} & \textbf{问题}\\ \midrule \end{tabularx}}{\ifdefstring{\ParallelMode}{right}{\phantomsection\label{notes-table}\hypertarget{notes-table}{}}{}\begin{tabularx}{\linewidth}{@{}>{\RaggedRight\arraybackslash}X>{\RaggedRight\arraybackslash}X@{}}\toprule \textbf{観察} & \textbf{問い}\\ \midrule \end{tabularx}}

\ParallelText{notes-table.row-1}{\begin{tabularx}{\linewidth}{@{}>{\RaggedRight\arraybackslash}X>{\RaggedRight\arraybackslash}X@{}}一张空长椅 & 换个时间会有人坐吗？\\ \end{tabularx}}{\begin{tabularx}{\linewidth}{@{}>{\RaggedRight\arraybackslash}X>{\RaggedRight\arraybackslash}X@{}}空いているベンチ & 別の時間には使われている？\\ \end{tabularx}}

\ParallelText{notes-table.row-2}{\begin{tabularx}{\linewidth}{@{}>{\RaggedRight\arraybackslash}X>{\RaggedRight\arraybackslash}X@{}}一扇新刷漆的门 & 它以前是什么颜色？\\ \end{tabularx}}{\begin{tabularx}{\linewidth}{@{}>{\RaggedRight\arraybackslash}X>{\RaggedRight\arraybackslash}X@{}}塗り直された扉 & 前は何色だった？\\ \end{tabularx}}

\ParallelText{notes-table.row-3}{\begin{tabularx}{\linewidth}{@{}>{\RaggedRight\arraybackslash}X>{\RaggedRight\arraybackslash}X@{}}窗户里传来音乐 & 走远一点，听起来会不同吗？\\ \bottomrule \end{tabularx}}{\begin{tabularx}{\linewidth}{@{}>{\RaggedRight\arraybackslash}X>{\RaggedRight\arraybackslash}X@{}}窓から聞こえる音楽 & 離れると聞こえ方は変わる？\\ \bottomrule \end{tabularx}}

\ParallelText{notes-table.caption}{这些虚构的笔记展示了观察与问题的区别，只是示例，并非对某个真实街区的调查结果。}{架空のメモで、観察と問いの違いを示しています。実在する地域についての調査結果ではありません。}\end{ParallelKeep}

\ParallelSection{return}{回来以后，再读一遍这段路}{歩いた道を読み返す}

\ParallelText{return-prose}{回来后，可以选一条笔记，展开成一小段文字。保留观察到的细节，说明它为什么对你重要，并标明还有哪些事情尚不清楚。事实性的问题可以稍后核实，但不必把每次散步都变成一个研究项目。有时候，认真留意过，本身就有价值。}{帰ったら、メモを一つ選んで短い段落にしてみます。観察した細部を残し、自分にとって気になった理由を書き、まだ分からないことも明示します。事実に関する問いは後で確かめられますが、散歩を毎回研究課題にする必要はありません。ただ注意を向けたことに価値がある場合もあります。}

\ParallelText{route-reference}{\ParallelReference{route}{回看路线部分}}{\ParallelReference{route}{道順の節に戻る}}
